\begin{definition}[The normalized Gaussian Riesz kernel $K_{\gamma,j}$]
Fix $n \geq 1$ and $j \in \{1,\ldots,n\}$.  Define, for $x, y \in \mathbb{R}^n$,
\[
  K_{\gamma,j}(x,y) = \frac{\sqrt2}{\pi^{(n+1)/2}}\,k_j^{\mathrm{raw}}(x,y),
  \qquad
  k_j^{\mathrm{raw}}(x,y)
    = \int_0^1 \frac{y_j - r x_j}{(1-r^2)^{n/2+1}\sqrt{-\log r}}
        \exp\Bigl(-\frac{|rx-y|^2}{1-r^2}\Bigr) dr ,
\]
the $r$-integral being taken over the open interval $(0,1)$ with respect to
Lebesgue measure.  This is the off-diagonal Lebesgue kernel of the $j$th
coordinate $R_{\gamma,j}$ of the tablet's energy-normalized Gaussian Riesz
transform, as displayed in the reference paper.  The displayed formula is the
whole of the definition: it is an explicit function of $(x,y)$ and depends on no
other tablet node, which is why this node imports only the shared preamble root.

\emph{Where the formula comes from.}  It is the time integral of the
$x$-derivative of the Lebesgue kernel, reparametrized by $r = e^{-s}$.  Writing
the Riesz transform as $(1/\sqrt\pi)\int_0^\infty D^x_j P_s\,ds/\sqrt s$ and
inserting the derivative identity
$D_j^x p_s(x,y) = \sqrt2\,r(y_j-rx_j)(1-r^2)^{-1}p_s(x,y)$ for the Lebesgue
kernel $p_s$ of the Ornstein--Uhlenbeck semigroup, the substitution $r = e^{-s}$, under
which $s = -\log r$, $ds = -dr/r$ and $s \in (0,\infty)$ corresponds to
$r \in (0,1)$, produces exactly the displayed integrand: the factor $r$ from the
derivative cancels the $1/r$ of $ds$, the $\pi^{-n/2}$ of $p_s$ combines with the
$1/\sqrt\pi$ into $\pi^{-(n+1)/2}$, the two powers of $1-r^2$ combine into
$(1-r^2)^{-n/2-1}$, and $s^{-1/2}$ becomes $(-\log r)^{-1/2}$.  This paragraph
records provenance only; the object defined here is the displayed formula, and
the identification of $K_{\gamma,j}$ with the operator is a separate assertion
proved elsewhere.

\emph{Scope of the definition.}  The reference paper displays the formula for
$x \neq y$, which is the only range in which it is used.  For $x = y$ the
$r$-integral need not converge, and the definition then takes the value the
Bochner integral convention assigns to a non-integrable integrand, namely $0$.
No statement built on this definition refers to the value on the diagonal.  The
integrand is also left unconstrained at the two endpoints $r = 0$ and $r = 1$,
which the open interval excludes; they form a Lebesgue-null set and do not affect
the integral.

\emph{The normalization is part of the definition, not a convention that can be
dropped.}  The function $k_j^{\mathrm{raw}}$ alone is the kernel of the
historical reference, and the prefactor $\sqrt2\,\pi^{-(n+1)/2}$ is what makes
$K_{\gamma,j}$ the kernel of the energy-normalized $R_{\gamma,j}$ rather than of
a constant multiple of it; in the same letters, the kernel of the historical
spectral operator $\mathcal{R}_j$ is $-2\pi^{-(n+1)/2}k_j^{\mathrm{raw}}$.  As
the paper states, a claim about dimension-free constants must specify these
factors, since an unnormalized kernel is a different family of operators.
\end{definition}
